\section{Predictions}

A theory earns its keep by what it risks. Table~\ref{tab:status} gathers the headline results with their status, derived by argument, measured on the substrate against a control, or still open.

\begin{table}[ht]
\centering
\small
\begin{tabular}{@{}lll@{}}
\toprule
result & sector & status \\
\midrule
spinors and the double cover on the 24 directions & matter & derived, measured \\
fermions by topology, spin-statistics & matter & derived \\
$g$-factor $= 2$ & matter & measured \\
mass as the chirality coupling, single chiral species & matter & measured \\
$\mathrm{SO}(10)$ unification, one generation per 16-spinor & force & derived \\
$\sin^2\theta_W = 3/8$, running to $0.231$ & force & derived \\
Standard-Model hypercharges from anomaly cancellation & force & derived \\
confinement, lattice quantum electrodynamics & force & measured \\
emergent Lorentz invariance, isotropic to fourth order & spacetime & measured, control \\
Newtonian $1/r$, Einstein as an equation of state & gravity & measured \\
Ryu-Takayanagi area law on the lower face & gravity & measured, control \\
de Sitter expansion, three space dimensions selected & cosmos & measured, control \\
baryogenesis meeting the three Sakharov conditions & cosmos & measured \\
$\mathrm{CHSH} = 2\sqrt{2}$ from a local base & quantum & measured, control \\
three generations carry distinct masses & matter & open conjecture \\
the coupling value, curved-bulk gravity, measurement & various & open \\
\bottomrule
\end{tabular}
\caption{The status of each headline claim. A result counts only where its control could have failed.}
\label{tab:status}
\end{table}

A few of these are genuine predictions, sharp enough to be wrong. The weak mixing angle is fixed at $3/8$ at unification and must run to the measured value at low energy, which it does. Space must have exactly three dimensions, because only there do orbits stay closed, a retrodiction the model cannot wriggle out of. And the discreteness leaves a faint fingerprint, a violation of Lorentz invariance that grows as the square of the energy over the Planck scale, far too small to see today but falsifiable in principle, and already enough to exclude a flat lattice while the hyperbolic substrate passes. The boldest open bet is the three generations, where the geometry forces the threefold structure but not yet the three masses, and a clean derivation of the splitting, or a clean proof that the structure is only one generation in disguise, would settle the program's central wager either way.

Three further bets follow from the committed structure, and each could fail. The matter spectrum is closed, exactly three generations of one $\mathrm{SO}(10)$ spinor apiece and nothing besides, so a fourth generation or a supersymmetric partner found at any collider would break it outright. Dark matter is not a particle here but a feature of the curved substrate, a geometric stand-in for the missing pull, so the long null result of the direct searches is the expected one, and a clean detection of a dark-matter particle would count against the model. And the same $\mathrm{SO}(10)$ that unifies the forces lets the proton decay, the generic fate of a grand-unified theory, at a rate the still-open coupling does not yet fix, so a measured proton lifetime would pin that last number rather than merely test the idea. None of these asks for new machinery. Each is forced by the structure already on the table, and each is sharp enough to lose.
